As métricas apresentadas neste trabalho utilizam um procedimento comum de tratamento dos registros de interação (\textit{traces}). Primeiramente, os \textit{traces} são agrupados por usuário, de forma que cada conjunto de registros corresponda a um único aluno. Em seguida, os \textit{timestamps} são convertidos em objetos temporais e ordenados cronologicamente, permitindo a reconstrução da sequência real de eventos.

Após a ordenação, procede-se ao seccionamento temporal, que consiste em analisar os intervalos entre \textit{traces} consecutivos. Para as métricas de tempo médio de sessão, engajamento por câmera e engajamento por microfone, esses intervalos são utilizados para determinar durações totais de sessão; para as métricas de distração e fragmentação do foco, eles permitem identificar transições entre estados (como ``em foco'', ``distração'' ou ``engajamento multimodal''); e para as métricas de taxa e frequência de distração, os intervalos são empregados no cálculo de taxas e frequências de ocorrência, assim como na extração de períodos contínuos associados a diferentes condições de engajamento.

Dessa forma, todas as análises descritas nas subseções posteriores fazem uso desta rotina padronizada composta por: (i) agrupamento por aluno; (ii) ordenação temporal; e (iii) análise intervalar entre registros, variando apenas quanto aos critérios específicos de classificação e contagem definidos para cada métrica.
